\chapter{Indirizzi MAC, switch, ARP ed Ethernet}

% ══════════════════════════════════════════════════════════════
\section{Gli indirizzi MAC}
% ══════════════════════════════════════════════════════════════

\dfn{Indirizzo MAC}{
    A livello di collegamento i dispositivi sono identificati dagli \textbf{indirizzi MAC} (o \textbf{indirizzi fisici}).
    \begin{itemize}
      \item Ogni interfaccia di rete ha un proprio indirizzo MAC (pensato per essere fisso, anche se si può cambiare).
      \item Ogni produttore ha un proprio insieme di indirizzi MAC (i primi 3 byte identificano il produttore).
      \item Un indirizzo MAC è lungo \textbf{6 byte} ($2^{48}$ indirizzi possibili) e si scrive di solito in esadecimale:
            \texttt{1A:23:F9:CD:06:9B} oppure \texttt{1A-23-F9-CD-06-9B}.
    \end{itemize}
}

\begin{table}[H]
\centering
\begin{tabular}{p{3cm}p{5.5cm}p{5.5cm}}
\toprule
 & \textbf{Indirizzo IP} & \textbf{Indirizzo MAC} \\
\midrule
Livello & rete (3) & collegamento (2) \\
Lunghezza & 32 bit (IPv4) & 48 bit \\
Notazione & decimale puntata \texttt{192.168.1.10} & esadecimale \texttt{1A:23:F9:CD:06:9B} \\
Portata & \textbf{globale} (fuori dalle LAN) & \textbf{locale} (solo nella LAN) \\
Assegnazione & dipende dalla rete in cui ci si trova & dal produttore della scheda \\
Analogia & l'indirizzo di casa (cambia se traslochi) & il codice fiscale (ti segue ovunque) \\
\bottomrule
\end{tabular}
\caption{Indirizzi IP e MAC a confronto.}
\end{table}

\warning{Gli indirizzi MAC sono locali}{
    Tutte le interfacce hanno un indirizzo MAC, ma questo è usato \textbf{solo all'interno di una LAN}: un frame non esce dalla propria rete locale.
    Gli indirizzi IP, invece, fuori dalle LAN sono globali. Mentre il datagramma IP viaggia dalla sorgente alla destinazione mantenendo gli stessi
    indirizzi IP, gli indirizzi MAC del frame \textbf{cambiano a ogni hop}.
}

\subsection{Come lavorano gli indirizzi MAC}

In una LAN, due interfacce A e B comunicano così:

\begin{itemize}
    \item A inserisce nel frame l'indirizzo MAC di B e lo trasmette;
    \item B riceve il frame e confronta il proprio indirizzo MAC con il MAC di destinazione del frame;
    \item se coincidono il frame viene \textbf{accettato}, altrimenti viene \textbf{scartato} (e il resto della pila non viene nemmeno coinvolto).
\end{itemize}

\dfn{Indirizzo di broadcast}{
    Si possono anche inviare messaggi \textbf{broadcast}, accettati da tutti indipendentemente dal MAC. Nelle LAN con indirizzi da 6 byte
    (come Ethernet e 802.11) l'indirizzo broadcast è una stringa di \textbf{48 uni}: \texttt{FF-FF-FF-FF-FF-FF}.
}

In una LAN è del tutto possibile (anzi frequente) che un'interfaccia riceva frame destinati a un'altra: il ruolo dell'indirizzo MAC è
\textbf{filtrare} i frame non destinati a sé, senza disturbare l'host.

% ══════════════════════════════════════════════════════════════
\section{Indirizzi e livelli}
% ══════════════════════════════════════════════════════════════

\begin{figure}[H]
\centering
\begin{tikzpicture}[every node/.style={font=\footnotesize\sffamily}]
  \node[lay app, minimum width=3cm] (a) at (0,2.4) {Applicazione};
  \node[lay tra, minimum width=3cm] (t) at (0,1.8) {Trasporto};
  \node[lay net, minimum width=3cm] (n) at (0,1.2) {Rete};
  \node[lay link, minimum width=3cm] (l) at (0,0.6) {Collegamento};
  \node[lay phy, minimum width=3cm] (p) at (0,0) {Fisico};
  \draw[decorate, decoration={brace, amplitude=6pt, mirror}] (-1.7,2.7) -- (-1.7,0.95) node[midway, left=8pt, text=netgreen!50!black, align=right] {indirizzi IP\\(e nomi, porte)};
  \node[anchor=west, text=netred] (mac) at (2.2,0.6) {indirizzi MAC};
  \draw[-{Stealth}, netred] (mac) -- (l.east);
  \node[anchor=west, text=primary, text width=5cm] at (2.2,1.35) {\textbf{ARP} fa da ponte tra i due tipi di indirizzi};
\end{tikzpicture}
\caption{Nella pila IP e MAC convivono: ARP collega gli uni agli altri (una differenza rispetto al modello ISO/OSI puro).}
\end{figure}

% ══════════════════════════════════════════════════════════════
\section{Gli switch}
% ══════════════════════════════════════════════════════════════

\dfn{Switch}{
    Gli \textbf{switch} sono l'equivalente dei router al livello di collegamento:
    \begin{itemize}
      \item non implementano alcun algoritmo di routing;
      \item usano \textbf{solo indirizzi MAC}, non indirizzi IP.
    \end{itemize}
    Il loro compito è ricevere i frame in arrivo e \textbf{inoltrarli sui collegamenti in uscita} giusti.
}

\begin{itemize}
    \item Lo switch è \textbf{trasparente} a host e router della subnet: nessuno gli indirizza frame, e nessuno si accorge che c'è.
    \item Uno switch ha dei \textbf{buffer} sulle interfacce.
    \item Gli switch hanno \textbf{tabelle di inoltro} che associano indirizzi MAC a interfacce.
    \item La tabella si aggiorna \textbf{automaticamente e dinamicamente} man mano che si scoprono nuovi dispositivi (\textbf{self-learning}).
\end{itemize}

\begin{figure}[H]
\centering
\begin{tikzpicture}[every node/.style={font=\scriptsize\sffamily}]
  \node (sw) at (0,0) {\icon[1.4cm]{switch}};
  \node (a) at (-3.2,1.2) {\icon[0.8cm]{pc}}; \node[left=0pt of a, font=\scriptsize\ttfamily] {A: 1A-23-F9-CD-06-9B};
  \node (b) at (-3.2,-1.2) {\icon[0.9cm]{laptop}}; \node[left=0pt of b, font=\scriptsize\ttfamily] {B: 5C-66-AB-90-75-B1};
  \node (c) at (3.2,1.2) {\icon[0.8cm]{pc}}; \node[right=0pt of c, font=\scriptsize\ttfamily] {C: 49-BD-D2-C7-56-2A};
  \node (d) at (3.2,-1.2) {\icon[0.9cm]{laptop}}; \node[right=0pt of d, font=\scriptsize\ttfamily] {D: 88-B2-2F-54-1A-0F};
  \draw[wired] (a) -- node[above, sloped] {porta 1} (sw); \draw[wired] (b) -- node[below, sloped] {porta 2} (sw);
  \draw[wired] (c) -- node[above, sloped] {porta 3} (sw); \draw[wired] (d) -- node[below, sloped] {porta 4} (sw);
  \node[draw, fill=llink!40, anchor=north, font=\scriptsize\sffamily] at (0,-1.9) {
    \begin{tabular}{llc}
    \textbf{MAC} & \textbf{interfaccia} & \textbf{TTL} \\
    1A-23-F9-CD-06-9B & 1 & 60 s \\
    88-B2-2F-54-1A-0F & 4 & 60 s \\
    \end{tabular}};
\end{tikzpicture}
\caption{Uno switch e la sua tabella, costruita osservando i frame che arrivano.}
\end{figure}

\subsection{Il self-learning}

\begin{enumerate}
    \item Quando arriva un frame sulla porta $x$, lo switch \textbf{impara} che il MAC \textbf{sorgente} si trova dietro la porta $x$ e lo registra
          nella tabella (con un tempo di vita).
    \item Poi cerca nella tabella il MAC di \textbf{destinazione}:
    \begin{itemize}
      \item se lo trova e la porta è diversa da $x$, \textbf{inoltra} il frame solo su quella porta;
      \item se la porta coincide con $x$, \textbf{scarta} il frame (il destinatario è dallo stesso lato: lo ha già ricevuto);
      \item se non lo trova, \textbf{inonda} (\emph{flooding}): invia il frame su tutte le porte tranne $x$.
    \end{itemize}
\end{enumerate}

\begin{esempio}{Lo switch impara}
Tabella vuota. A (porta 1) invia un frame a D.
\begin{itemize}
    \item Lo switch impara ``A è sulla porta 1''. Non conosce D: inonda le porte 2, 3 e 4.
    \item D risponde ad A. Lo switch impara ``D è sulla porta 4''; conosce A: inoltra solo sulla porta 1.
    \item Da ora in poi i frame tra A e D viaggiano solo tra le porte 1 e 4.
\end{itemize}
\end{esempio}

% ══════════════════════════════════════════════════════════════
\section{La LAN commutata}
% ══════════════════════════════════════════════════════════════

\begin{figure}[H]
\centering
\begin{minipage}{0.5\linewidth}\centering
  \includegraphics[width=0.95\linewidth]{cap20/switched_lan.png}
\end{minipage}\hfill
\begin{minipage}{0.46\linewidth}
È comune che le LAN usino uno o più switch per collegare molti dispositivi locali (nell'esempio del Kurose, i dipartimenti di un
istituto collegati da switch, con un router verso Internet e i server Web e di posta).
\end{minipage}
\caption{Una rete di istituto con switch organizzati ad albero (da Kurose).}
\end{figure}

A differenza dei router, gli switch sono \textbf{veloci} e \textbf{plug-and-play}:

\begin{itemize}
    \item non c'è alcun algoritmo di routing;
    \item si considerano solo i primi \textbf{2 livelli} della pila.
\end{itemize}

D'altra parte le LAN commutate sono \textbf{limitate in dimensione} e devono essere \textbf{strutturate ad albero}:

\begin{itemize}
    \item gli indirizzi MAC sono difficili da raggruppare (non hanno prefissi legati alla posizione): le tabelle degli switch possono crescere rapidamente;
    \item non c'è routing: i \textbf{cicli} sono difficili da evitare (un frame inondato in un ciclo girerebbe per sempre).
\end{itemize}

\begin{table}[H]
\centering
\begin{tabular}{p{3.5cm}p{5.2cm}p{5.2cm}}
\toprule
 & \textbf{Switch} & \textbf{Router} \\
\midrule
Livello & collegamento (2) & rete (3) \\
Indirizzi & MAC & IP \\
Tabelle & self-learning & algoritmi di routing \\
Configurazione & plug-and-play & va configurato \\
Topologia & albero (niente cicli) & qualsiasi, anche con cicli \\
Traffico broadcast & lo inoltra & lo blocca (separa i domini) \\
\bottomrule
\end{tabular}
\caption{Switch e router a confronto.}
\end{table}

% ══════════════════════════════════════════════════════════════
\section{Il protocollo ARP}
% ══════════════════════════════════════════════════════════════

I protocolli dei livelli superiori lavorano con gli indirizzi IP, mentre il frame ha bisogno dell'indirizzo MAC del destinatario. Serve quindi
tradurre gli indirizzi IP in indirizzi MAC.

\dfn{ARP (Address Resolution Protocol)}{
    L'\textbf{ARP} gestisce la conversione tra indirizzi IP e indirizzi MAC.
    \begin{itemize}
      \item Ogni interfaccia ha un modulo ARP con una \textbf{tabella ARP}, che associa ogni IP della LAN a un indirizzo MAC, con un
            \textbf{TTL} (\emph{time to live}) scaduto il quale la voce viene cancellata (tipicamente 5--20 minuti).
      \item Poiché gli indirizzi MAC sono locali, anche ARP funziona \textbf{solo nelle reti locali}.
    \end{itemize}
}

\begin{table}[H]
\centering
\begin{tabular}{lll}
\toprule
\textbf{Indirizzo IP} & \textbf{Indirizzo MAC} & \textbf{TTL} \\
\midrule
222.222.222.221 & 88-B2-2F-54-1A-0F & 13:45:00 \\
222.222.222.223 & 5C-66-AB-90-75-B1 & 13:52:00 \\
\bottomrule
\end{tabular}
\caption{Una tabella ARP: la voce scade all'orario indicato.}
\end{table}

\info{ARP e DNS}{
    ARP fa per la LAN quello che il DNS fa per Internet: il DNS traduce \textbf{nomi} in indirizzi IP (in tutta Internet), ARP traduce
    \textbf{indirizzi IP} in indirizzi MAC (solo nella propria LAN). Sulle proprie macchine si può vedere la tabella con \texttt{ip neigh}
    (o \texttt{arp -a}).
}

\subsection{Un esempio di ARP}

\begin{figure}[H]
\centering
\begin{minipage}{0.44\linewidth}\centering
  \includegraphics[width=0.95\linewidth]{cap20/arp_lan.png}
\end{minipage}\hfill
\begin{minipage}{0.52\linewidth}
L'host C (\texttt{222.222.222.220}) vuole inviare messaggi all'host A (\texttt{222.222.222.222}). Per farlo deve conoscere anche il MAC di A.
\begin{enumerate}
  \item Se A \textbf{non è nella tabella}, prima di inviare il messaggio C invia in \textbf{broadcast} a tutti i dispositivi della rete un
        pacchetto \textbf{ARP request}: ``chi ha l'IP \texttt{222.222.222.222}? Lo dica a \texttt{222.222.222.220}''.
  \item Tutti i nodi ricevono il pacchetto, ma \textbf{solo} chi ha l'IP cercato risponde, con un messaggio \textbf{ARP reply} diretto
        (\emph{unicast}) che contiene il proprio MAC.
  \item C inserisce la coppia IP--MAC nella propria tabella e invia il frame.
\end{enumerate}
\end{minipage}
\caption{ARP in una LAN: tre host e un router collegati da uno switch (da Kurose).}
\end{figure}

\begin{figure}[H]
\centering
\begin{tikzpicture}[yscale=0.85, every node/.style={font=\scriptsize\sffamily}]
  \timeline{c}{0}{C}{3.6}
  \timeline{b}{4}{B}{3.6}
  \timeline{a}{8}{A}{3.6}
  \draw[msg, netred] (0,-0.5) -- node[above, sloped, pos=0.5] {ARP request (dest. FF-FF-FF-FF-FF-FF): chi ha 222.222.222.222?} (8,-0.5);
  \fill[netred] (4,-0.5) circle (2.5pt); \node[anchor=west, text=netred] at (4.1,-0.85) {riceve e ignora};
  \node[anchor=west] at (8.2,-0.5) {``sono io!''};
  \draw[msg, netgreen!60!black] (8,-1.7) -- node[above, pos=0.5] {ARP reply (unicast): 222.222.222.222 è 49-BD-D2-C7-56-2A} (0,-1.7);
  \node[anchor=east, text width=2.6cm, align=right] at (-0.2,-1.9) {aggiorna la tabella ARP};
  \draw[msg, netblue, very thick] (0,-2.8) -- node[above] {frame IP con MAC di destinazione di A} (8,-2.8);
\end{tikzpicture}
\caption{Richiesta in broadcast, risposta in unicast.}
\end{figure}

\warning{ARP poisoning}{
    ARP non prevede alcuna autenticazione: un host \textbf{malevolo} può intercettare la richiesta e rispondere con una risposta
    \textbf{falsificata}, associando l'IP cercato al proprio MAC (\textbf{ARP poisoning}). I frame destinati alla vittima arrivano così
    all'attaccante, che può leggerli o modificarli (\emph{man in the middle}). È uno degli attacchi più comuni nelle LAN.
}

\subsection{ARP e il gateway}

Che succede se l'IP di destinazione è \textbf{fuori dalla rete}?

\begin{figure}[H]
    \centering
    \includegraphics[width=0.7\linewidth]{cap20/arp_gateway.png}
    \caption{Due subnet (111.111.111.x e 222.222.222.x) collegate da un router con due interfacce, ciascuna con il proprio IP, MAC e tabella ARP.}
\end{figure}

\begin{itemize}
    \item Il router che collega le due reti deve avere almeno \textbf{2 interfacce} (2 IP, 2 MAC e 2 tabelle ARP), una dentro ciascuna subnet.
    \item I frame diretti fuori dalla subnet vengono inviati alla \textbf{prima interfaccia del router}: l'host usa ARP per il MAC del
          \textbf{router} (il gateway), non per quello del destinatario finale.
    \item Il router estrae il datagramma, lo sposta sulla seconda interfaccia e, con la seconda tabella ARP, lo invia all'host giusto in un
          \textbf{nuovo frame}.
\end{itemize}

\begin{esempio}{Da A (111.111.111.111) a B (222.222.222.222)}
\begin{enumerate}
    \item A capisce, confrontando indirizzo e maschera, che B è in un'altra subnet: deve passare dal gateway \texttt{111.111.111.110}.
    \item A usa ARP per trovare il MAC dell'interfaccia del router \texttt{111.111.111.110} (\texttt{E6-E9-00-17-BB-4B}).
    \item A crea un datagramma IP con \textbf{IP sorgente A} e \textbf{IP destinazione B}, e lo incapsula in un frame con
          \textbf{MAC destinazione = MAC del router}.
    \item Il router riceve il frame, estrae il datagramma, vede che B è sulla subnet \texttt{222.222.222.0/24}, usa ARP su quell'interfaccia per
          trovare il MAC di B e crea un nuovo frame con \textbf{MAC sorgente = sua seconda interfaccia} e \textbf{MAC destinazione = B}.
\end{enumerate}
\begin{center}
\begin{tabular}{lcccc}
\toprule
\textbf{Tratta} & \textbf{MAC sorgente} & \textbf{MAC destinazione} & \textbf{IP sorgente} & \textbf{IP destinazione} \\
\midrule
A $\rightarrow$ router & A & router (if. 1) & A & B \\
router $\rightarrow$ B & router (if. 2) & B & A & B \\
\bottomrule
\end{tabular}
\end{center}
Gli indirizzi IP restano gli stessi da un estremo all'altro; gli indirizzi MAC cambiano a ogni collegamento.
\end{esempio}

% ══════════════════════════════════════════════════════════════
\section{Il frame Ethernet}
% ══════════════════════════════════════════════════════════════

Ethernet è la tecnologia dominante per le LAN cablate. Nata negli anni '70 come bus condiviso con CSMA/CD, oggi è quasi sempre una
topologia a stella con uno switch al centro e collegamenti full-duplex (dove le collisioni non avvengono più).

\begin{figure}[H]
\centering
\begin{tikzpicture}[every node/.style={font=\scriptsize\sffamily}]
  \node[anchor=east, font=\footnotesize\bfseries\sffamily] at (-0.2,0) {Ethernet (DIX)};
  \node[field, fill=lphy, minimum width=1.6cm] (pr) at (0.8,0) {preambolo};
  \node[field, fill=llink, minimum width=1.9cm, right=0pt of pr] (dst) {indirizzo dest.};
  \node[field, fill=llink, minimum width=1.9cm, right=0pt of dst] (src) {indirizzo sorg.};
  \node[field, fill=lnet!50, minimum width=0.9cm, right=0pt of src] (ty) {tipo};
  \node[field, fill=cloudbg, minimum width=3.2cm, right=0pt of ty] (da) {dati};
  \node[field, fill=cloudbg!50, minimum width=0.9cm, right=0pt of da] (pad) {pad};
  \node[field, fill=netred!25, minimum width=1.3cm, right=0pt of pad] (crc) {CRC};
  \foreach \n/\b in {pr/8, dst/6, src/6, ty/2, da/{0--1500}, pad/{0--46}, crc/4} \node[above=0pt of \n] {\b};
  \node[anchor=east] at (-0.2,0.45) {byte};
  \begin{scope}[yshift=-1.1cm]
    \node[anchor=east, font=\footnotesize\bfseries\sffamily] at (-0.2,0) {IEEE 802.3};
    \node[field, fill=lphy, minimum width=1.3cm] (pr2) at (0.65,0) {preambolo};
    \node[field, fill=lphy!60, minimum width=0.3cm, right=0pt of pr2] (sof) {\rotatebox{90}{\tiny SOF}};
    \node[field, fill=llink, minimum width=1.9cm, right=0pt of sof] (dst2) {indirizzo dest.};
    \node[field, fill=llink, minimum width=1.9cm, right=0pt of dst2] (src2) {indirizzo sorg.};
    \node[field, fill=lnet!50, minimum width=0.9cm, right=0pt of src2] (ln) {lungh.};
    \node[field, fill=cloudbg, minimum width=3.2cm, right=0pt of ln] {dati};
    \node[field, fill=cloudbg!50, minimum width=0.9cm] at (9.95,0) {pad};
    \node[field, fill=netred!25, minimum width=1.3cm] at (11.05,0) {CRC};
  \end{scope}
\end{tikzpicture}
\caption{Il frame Ethernet nelle due varianti (DIX e IEEE 802.3).}
\end{figure}

\begin{itemize}
    \item \textbf{Preambolo} (8 byte): un blocco ``sveglia'' di bit alternati (\texttt{10101010} \dots, con l'ultimo byte \texttt{10101011})
          usato per \textbf{sincronizzare gli orologi} degli adattatori di sorgente e destinazione.
    \item \textbf{Indirizzo di destinazione} (6 byte): il MAC dell'adattatore di destinazione.
    \item \textbf{Indirizzo sorgente} (6 byte): il MAC dell'adattatore che trasmette il frame nella LAN.
    \item \textbf{Tipo} (2 byte): il protocollo di livello di rete contenuto nel frame. Non c'è solo IP (\texttt{0x0800}): per esempio i pacchetti
          ARP hanno un tipo specifico, \texttt{0x0806}. Nella variante 802.3 al suo posto c'è la lunghezza.
    \item \textbf{Dati} (da 46 a 1500 byte): il datagramma IP. Il massimo è la \textbf{MTU} di Ethernet: se il datagramma è più grande, va frammentato.
    \item \textbf{Pad}: riempimento per raggiungere la lunghezza minima dei dati (46 byte), necessaria perché il CSMA/CD rilevi le collisioni
          (tempo di trasmissione di almeno $2\tau$).
    \item \textbf{CRC} (4 byte): il CRC-32 calcolato sul frame, per rilevare gli errori.
\end{itemize}

\info{Un servizio senza connessione e inaffidabile}{
    Ethernet è \textbf{senza connessione} (nessun handshake con l'adattatore di destinazione) e \textbf{inaffidabile}: il ricevente non invia né ACK
    né NAK. Un frame che non supera il controllo CRC viene semplicemente scartato; se i dati servono davvero, sarà TCP (o l'applicazione) a
    accorgersi della mancanza e a ritrasmettere.
}

% ══════════════════════════════════════════════════════════════
\section{Riepilogo}
% ══════════════════════════════════════════════════════════════

\riepilogo{
\begin{itemize}
    \item L'\textbf{indirizzo MAC} (48 bit, esadecimale) identifica un'interfaccia ed è \textbf{locale}; broadcast = \texttt{FF-FF-FF-FF-FF-FF}.
    \item Gli \textbf{switch} inoltrano frame in base al MAC, con tabelle \textbf{self-learning}; sono veloci e plug-and-play ma le LAN commutate
          devono essere ad albero.
    \item \textbf{ARP} traduce IP in MAC: richiesta in broadcast, risposta in unicast, tabella con TTL; vulnerabile all'\textbf{ARP poisoning}.
    \item Verso un'altra subnet il frame va al MAC del \textbf{router}: gli IP restano uguali, i MAC cambiano a ogni hop.
    \item \textbf{Frame Ethernet}: preambolo (8), destinazione (6), sorgente (6), tipo (2), dati (46--1500, con pad), CRC (4).
\end{itemize}
}
